% Options for packages loaded elsewhere
% Options for packages loaded elsewhere
\PassOptionsToPackage{unicode}{hyperref}
\PassOptionsToPackage{hyphens}{url}
\PassOptionsToPackage{dvipsnames,svgnames,x11names}{xcolor}
%
\documentclass[
  letterpaper,
  DIV=11,
  numbers=noendperiod]{scrartcl}
\usepackage{xcolor}
\usepackage{amsmath,amssymb}
\setcounter{secnumdepth}{-\maxdimen} % remove section numbering
\usepackage{iftex}
\ifPDFTeX
  \usepackage[T1]{fontenc}
  \usepackage[utf8]{inputenc}
  \usepackage{textcomp} % provide euro and other symbols
\else % if luatex or xetex
  \usepackage{unicode-math} % this also loads fontspec
  \defaultfontfeatures{Scale=MatchLowercase}
  \defaultfontfeatures[\rmfamily]{Ligatures=TeX,Scale=1}
\fi
\usepackage{lmodern}
\ifPDFTeX\else
  % xetex/luatex font selection
\fi
% Use upquote if available, for straight quotes in verbatim environments
\IfFileExists{upquote.sty}{\usepackage{upquote}}{}
\IfFileExists{microtype.sty}{% use microtype if available
  \usepackage[]{microtype}
  \UseMicrotypeSet[protrusion]{basicmath} % disable protrusion for tt fonts
}{}
\makeatletter
\@ifundefined{KOMAClassName}{% if non-KOMA class
  \IfFileExists{parskip.sty}{%
    \usepackage{parskip}
  }{% else
    \setlength{\parindent}{0pt}
    \setlength{\parskip}{6pt plus 2pt minus 1pt}}
}{% if KOMA class
  \KOMAoptions{parskip=half}}
\makeatother
% Make \paragraph and \subparagraph free-standing
\makeatletter
\ifx\paragraph\undefined\else
  \let\oldparagraph\paragraph
  \renewcommand{\paragraph}{
    \@ifstar
      \xxxParagraphStar
      \xxxParagraphNoStar
  }
  \newcommand{\xxxParagraphStar}[1]{\oldparagraph*{#1}\mbox{}}
  \newcommand{\xxxParagraphNoStar}[1]{\oldparagraph{#1}\mbox{}}
\fi
\ifx\subparagraph\undefined\else
  \let\oldsubparagraph\subparagraph
  \renewcommand{\subparagraph}{
    \@ifstar
      \xxxSubParagraphStar
      \xxxSubParagraphNoStar
  }
  \newcommand{\xxxSubParagraphStar}[1]{\oldsubparagraph*{#1}\mbox{}}
  \newcommand{\xxxSubParagraphNoStar}[1]{\oldsubparagraph{#1}\mbox{}}
\fi
\makeatother


\usepackage{longtable,booktabs,array}
\newcounter{none} % for unnumbered tables
\usepackage{calc} % for calculating minipage widths
% Correct order of tables after \paragraph or \subparagraph
\usepackage{etoolbox}
\makeatletter
\patchcmd\longtable{\par}{\if@noskipsec\mbox{}\fi\par}{}{}
\makeatother
% Allow footnotes in longtable head/foot
\IfFileExists{footnotehyper.sty}{\usepackage{footnotehyper}}{\usepackage{footnote}}
\makesavenoteenv{longtable}
\usepackage{graphicx}
\makeatletter
\newsavebox\pandoc@box
\newcommand*\pandocbounded[1]{% scales image to fit in text height/width
  \sbox\pandoc@box{#1}%
  \Gscale@div\@tempa{\textheight}{\dimexpr\ht\pandoc@box+\dp\pandoc@box\relax}%
  \Gscale@div\@tempb{\linewidth}{\wd\pandoc@box}%
  \ifdim\@tempb\p@<\@tempa\p@\let\@tempa\@tempb\fi% select the smaller of both
  \ifdim\@tempa\p@<\p@\scalebox{\@tempa}{\usebox\pandoc@box}%
  \else\usebox{\pandoc@box}%
  \fi%
}
% Set default figure placement to htbp
\def\fps@figure{htbp}
\makeatother





\setlength{\emergencystretch}{3em} % prevent overfull lines

\providecommand{\tightlist}{%
  \setlength{\itemsep}{0pt}\setlength{\parskip}{0pt}}



 


\KOMAoption{captions}{tableheading}
\makeatletter
\@ifpackageloaded{caption}{}{\usepackage{caption}}
\AtBeginDocument{%
\ifdefined\contentsname
  \renewcommand*\contentsname{Table of contents}
\else
  \newcommand\contentsname{Table of contents}
\fi
\ifdefined\listfigurename
  \renewcommand*\listfigurename{List of Figures}
\else
  \newcommand\listfigurename{List of Figures}
\fi
\ifdefined\listtablename
  \renewcommand*\listtablename{List of Tables}
\else
  \newcommand\listtablename{List of Tables}
\fi
\ifdefined\figurename
  \renewcommand*\figurename{Figure}
\else
  \newcommand\figurename{Figure}
\fi
\ifdefined\tablename
  \renewcommand*\tablename{Table}
\else
  \newcommand\tablename{Table}
\fi
}
\@ifpackageloaded{float}{}{\usepackage{float}}
\floatstyle{ruled}
\@ifundefined{c@chapter}{\newfloat{codelisting}{h}{lop}}{\newfloat{codelisting}{h}{lop}[chapter]}
\floatname{codelisting}{Listing}
\newcommand*\listoflistings{\listof{codelisting}{List of Listings}}
\makeatother
\makeatletter
\makeatother
\makeatletter
\@ifpackageloaded{caption}{}{\usepackage{caption}}
\@ifpackageloaded{subcaption}{}{\usepackage{subcaption}}
\makeatother
\usepackage{bookmark}
\IfFileExists{xurl.sty}{\usepackage{xurl}}{} % add URL line breaks if available
\urlstyle{same}
\hypersetup{
  pdftitle={FC02 --- Article 10 --- System Dynamics},
  colorlinks=true,
  linkcolor={blue},
  filecolor={Maroon},
  citecolor={Blue},
  urlcolor={Blue},
  pdfcreator={LaTeX via pandoc}}


\title{FC02 --- Article 10 --- System Dynamics}
\usepackage{etoolbox}
\makeatletter
\providecommand{\subtitle}[1]{% add subtitle to \maketitle
  \apptocmd{\@title}{\par {\large #1 \par}}{}{}
}
\makeatother
\subtitle{From Description to Model}
\author{}
\date{}
\begin{document}
\maketitle


\subsection{A formal definition of how the immune system
evolves}\label{a-formal-definition-of-how-the-immune-system-evolves}

Up to this point, we have been describing the system through patterns
and observations.

What follows is where that description becomes a model.

\begin{center}\rule{0.5\linewidth}{0.5pt}\end{center}

We have described the system.

Now we commit to modeling it.

If the immune system is a state-based system, one question becomes
unavoidable:

\textbf{How does the state evolve?}

A system is defined by how it evolves.\\
If we cannot describe that evolution, we do not yet have a model.

\begin{center}\rule{0.5\linewidth}{0.5pt}\end{center}

\subsection{1. System Definition}\label{system-definition}

We now make a formal commitment.

The immune system is a \textbf{partially observable, nonlinear,
distributed dynamical system with feedback}.

This is not a descriptive choice. It is a constraint imposed by the
system.

From this, one rule follows:

\begin{quote}
The future state of the system is determined by its current state and
the inputs acting on it.
\end{quote}

\begin{center}\rule{0.5\linewidth}{0.5pt}\end{center}

\subsection{2. State Representation}\label{state-representation}

We begin by defining the internal state of the system explicitly.

\[
x(t) \in \mathbb{R}^n
\]

Here, \(x(t)\) represents the \textbf{true internal state} of the immune
system at time \(t\).

This is not directly observable --- it is what we must infer.

\begin{center}\rule{0.5\linewidth}{0.5pt}\end{center}

\subsubsection{Explicit State Structure}\label{explicit-state-structure}

\[
x(t) =
\begin{bmatrix}
x_1(t) \\
x_2(t) \\
\vdots \\
x_n(t)
\end{bmatrix}
\]

Each component represents a dimension of the system:

\begin{itemize}
\tightlist
\item
  immune cell populations\\
\item
  activation states\\
\item
  cytokine environments\\
\item
  tissue condition\\
\item
  regulatory signals
\end{itemize}

The choice of components defines the resolution of the model.

\begin{center}\rule{0.5\linewidth}{0.5pt}\end{center}

\subsubsection{Partitioned State}\label{partitioned-state}

\[
x(t) =
\begin{bmatrix}
x_{\text{immune}}(t) \\
x_{\text{tissue}}(t) \\
x_{\text{regulatory}}(t)
\end{bmatrix}
\]

This reflects the distributed nature of the system.

\begin{center}\rule{0.5\linewidth}{0.5pt}\end{center}

\subsection{3. State Evolution}\label{state-evolution}

We now define how the state changes over time.

\[
x(t+1) = F\big(x(t), u(t)\big)
\]

The next state depends on the current state and what acts on it.

Where:

\begin{itemize}
\tightlist
\item
  \(F\) represents internal dynamics\\
\item
  \(u(t)\) represents inputs (therapy, infection, environment)
\end{itemize}

\begin{center}\rule{0.5\linewidth}{0.5pt}\end{center}

\subsubsection{Component-wise View}\label{component-wise-view}

\[
x_i(t+1) = f_i\big(x_1(t), \dots, x_n(t), u(t)\big)
\]

Each component evolves through interactions with others.

There is no central controller --- behavior emerges from interactions.

\begin{center}\rule{0.5\linewidth}{0.5pt}\end{center}

\subsection{3A. Minimal Example}\label{a.-minimal-example}

Consider a simple system:

\begin{itemize}
\tightlist
\item
  \(x(t)\) = inflammation level\\
\item
  \(u(t)\) = anti-inflammatory input
\end{itemize}

\begin{center}\rule{0.5\linewidth}{0.5pt}\end{center}

\subsubsection{Low initial state}\label{low-initial-state}

\begin{itemize}
\tightlist
\item
  \(x(0) = 0.2\) → decays → stabilizes
\end{itemize}

\begin{center}\rule{0.5\linewidth}{0.5pt}\end{center}

\subsubsection{Higher initial state}\label{higher-initial-state}

\begin{itemize}
\tightlist
\item
  \(x(0) = 0.6\) → amplifies → diverges
\end{itemize}

\begin{center}\rule{0.5\linewidth}{0.5pt}\end{center}

\subsubsection{Implication}\label{implication}

\begin{quote}
The trajectory depends on the state, not just the intervention.
\end{quote}

\begin{center}\rule{0.5\linewidth}{0.5pt}\end{center}

\subsection{3B. Basins of Attraction}\label{b.-basins-of-attraction}

The system evolves within a space of possible states.

Some states act as attractors.

\[
x^* = F(x^*, u^*)
\]

States near these points converge toward them.

\begin{center}\rule{0.5\linewidth}{0.5pt}\end{center}

\subsubsection{Interpretation}\label{interpretation}

\begin{quote}
Disease corresponds to being confined within a basin of attraction.
\end{quote}

\begin{center}\rule{0.5\linewidth}{0.5pt}\end{center}

\subsubsection{Therapy Reinterpreted}\label{therapy-reinterpreted}

\begin{itemize}
\tightlist
\item
  within basin → temporary effect\\
\item
  crossing basin → qualitative change
\end{itemize}

\begin{center}\rule{0.5\linewidth}{0.5pt}\end{center}

\subsection{3C. Energy Landscape
(Intuition)}\label{c.-energy-landscape-intuition}

We can interpret stability using a function:

\[
V(x)
\]

Lower values correspond to more stable states.

\begin{center}\rule{0.5\linewidth}{0.5pt}\end{center}

\subsubsection{Approximate dynamics}\label{approximate-dynamics}

\[
x(t+1) \approx x(t) - \nabla V(x(t))
\]

This is an approximation, not a general rule.

It helps visualize how systems settle into stable states.

\begin{center}\rule{0.5\linewidth}{0.5pt}\end{center}

\subsubsection{Implication}\label{implication-1}

\begin{quote}
Recovery requires crossing stability barriers, not just reversing
direction.
\end{quote}

\begin{center}\rule{0.5\linewidth}{0.5pt}\end{center}

\subsection{4. Properties of the
Dynamics}\label{properties-of-the-dynamics}

\subsubsection{Nonlinearity}\label{nonlinearity}

\[
F(x_1 + x_2, u) \neq F(x_1, u) + F(x_2, u)
\]

This allows thresholds and tipping behavior.

\begin{center}\rule{0.5\linewidth}{0.5pt}\end{center}

\subsubsection{State Dependence}\label{state-dependence}

\[
\frac{\partial F}{\partial u}(x_1) \neq \frac{\partial F}{\partial u}(x_2)
\]

The same input produces different outcomes in different states.

\begin{center}\rule{0.5\linewidth}{0.5pt}\end{center}

\subsubsection{Feedback}\label{feedback}

\[
x(t+2) = F(F(x(t), u(t)), u(t+1))
\]

The system evolves through feedback loops.

\begin{center}\rule{0.5\linewidth}{0.5pt}\end{center}

\subsubsection{Local Linear
Approximation}\label{local-linear-approximation}

\[
x(t+1) \approx A(x^*)x(t) + B(x^*)u(t)
\]

Valid only near a specific operating point.

\begin{center}\rule{0.5\linewidth}{0.5pt}\end{center}

\subsection{5. Observation}\label{observation}

We observe a projection of the true state. \[
y(t) = H(x(t))
\]

This means what we observe is a transformation of the underlying state,
not the state itself.

\begin{center}\rule{0.5\linewidth}{0.5pt}\end{center}

\subsubsection{Consequence}\label{consequence}

\[
x(t) \neq H^{-1}(y(t))
\]

The state cannot be uniquely reconstructed.

\begin{center}\rule{0.5\linewidth}{0.5pt}\end{center}

\subsection{6. Structure}\label{structure}

All previously introduced elements resolve into this structure:

\begin{itemize}
\tightlist
\item
  state\\
\item
  dynamics\\
\item
  observation\\
\item
  input
\end{itemize}

\begin{center}\rule{0.5\linewidth}{0.5pt}\end{center}

\subsection{7. Core Claim}\label{core-claim}

\begin{quote}
Disease is not a condition.\\
It is a trajectory of \(x(t)\) over time.
\end{quote}

And:

\begin{quote}
Clinical labels classify \(y(t)\), not \(x(t)\).
\end{quote}

\begin{center}\rule{0.5\linewidth}{0.5pt}\end{center}

\subsection{8. What This Model Does Not
Claim}\label{what-this-model-does-not-claim}

This is a model class, not a full specification.

It does not define:

\begin{itemize}
\tightlist
\item
  exact form of \(F\)\\
\item
  estimation methods\\
\item
  control strategies\\
\item
  disease-specific mappings
\end{itemize}

\begin{center}\rule{0.5\linewidth}{0.5pt}\end{center}

\subsection{Closing}\label{closing}

We no longer have a description.

We have a model class.

From this point forward:

\begin{quote}
Every claim must be consistent with this structure.
\end{quote}




\end{document}
